\chaptertitle{Three-Part Invention}

                                   Three-Part Invention

           Achilles (a Greek warrior, the fleetest of foot of all mortals) and a Tortoise are
           standing together on a dusty runway in the hot sun. Far down the runway, on a
           tall flagpole, there hangs a large rectangular flag. The flag is sold red, except
           where a thin ring-shaped holes has been cut out of it, through which one can see
           the sky.
\achilles What is that strange flag down at the other end of the track? It reminds me
     somehow of a print by my favourite artists M.C. Escher.
\tortoise That is Zeno’ s flag
\achilles Could it be that the hole in it resembles the holes in a Mobian strip Escher once
     drew? Something is wrong about the flag, I can tell.
\tortoise The ring which has been cut from it has the shape of the numeral for zero, which
     is Zeno´s favourite number.
\achilles The ring which hasn´t been invented yet! It will only be invented by a Hindu
     mathematician some millennia hence. And thus, Mr. T, mt argument proves that such a
     flag is impossible.
\tortoise Your argument is persuasive, Achilles, and I must agree that such a flag is indeed
     impossible. But it is beautiful anyway, is it not?
\achilles Oh, yes, there is no doubt of its beauty.
\tortoise I wonder if it´s beauty is related to it´s impossibility. I don´t know, I´ve never had
     the time to analyze Beauty. It´s a Capitalized Essence, and I never seem to have time for
     Capitalized Essences.
\achilles Speaking of Capitalized Essences, Mr. T, have you ever wondered about the
     Purpose of Life?
\tortoise Oh, heavens, no;
\achilles Haven’t you ever wondered why we are here, or who invented us?
\tortoise Oh, that is quite another matter. We are inventions of Zeno (as you will shortly
     see) and the reason we are here is to have a footrace.
ACHILLES::: A footrace? How outrageous! Me, the fleetest of foot of all mortals, versus you,
     the ploddingest of the plodders! There can be no point to such a race.
\tortoise You might give me a head start.
\achilles It would have to be a huge one.
\tortoise I don’t object.
\achilles But I will catch you, sooner or later – most likely sooner.
\tortoise Not if things go according to Zeno´s paradox, you won’t. Zeno is hoping to use
     our footrace to show that motion is impossible, you see. It is only in the mind that motion
     seems possible, according to Zeno. In truth, Motion Is Inherently Impossible. He proves
     it quite elegantly.




Three-Part Invention                                                                            37

Figure 10. Mobius strip by M.C.Escher (wood-engraving printed from four blocks, 1961)
\achilles Oh, yes, it comes back to me now: the famous Zen koan about Zen
       Master Zeno. As you say it is very simple indeed.
\tortoise Zen Koan? Zen Master? What do you mean?
\achilles It goes like this: Two monks were arguing about a flag. One said, “The
       flag is moving.” The other said, “The wind is moving.” The sixth patriarch, Zeno,
       happened to be passing by. He told them, “Not the wind, not the flag, mind is
       moving.”
\tortoise I am afraid you are a little befuddled, Achilles. Zeno is no Zen master, far
       from it. He is in fact, a Greek philosopher from the town of Elea (which lies halfway
       between points A and B). Centuries hence, he will be celebrated for his paradoxes of
       motion. In one of those paradoxes, this very footrace between you and me will play a
       central role.
\achilles I’m all confused. I remember vividly how I used to repeat over and over
       the names of the six patriarchs of Zen, and I always said, “The sixth patriarch is Zeno,
       The sixth patriarch is Zeno…” (Suddenly a soft warm breeze picks up.) Oh, look Mr.
       Tortoise – look at the flag waving! How I love to watch the ripples shimmer through
       it’s soft fabric. And the ring cut out of it is waving, too!



Three-Part Invention                                                                          38
\tortoise Don´t be silly. The flag is impossible, hence it can’t be waving. The wind is
     waving.

(At this moment, Zeno happens by.)
\zeno Hallo! Hulloo! What’s up? What’s new?
\achilles The flag is moving.
\tortoise The wind is moving.
\zeno O friends, Friends! Cease your argumentation! Arrest your vitriolics! Abandon your
      discord! For I shall resolve the issue for you forthwith. Ho! And on such a fine day.
\achilles This fellow must be playing the fool.
\tortoise No, wait, Achilles. Let us hear what he has to say. Oh Unknown Sir, do impart to
      us your thoughts on this matter.
\zeno Most willingly. Not thw ind, not the flag – neither one is moving, nor is anything moving
      at all. For I have discovered a great Theorem, which states; “Motion Is Inherently
      Impossible.” And from this Theorem follows an even greater Theorem – Zeno’s
      Theorem: “Motion Unexists.”
\achilles “Zeno’s Theorem”? Are you, sir, by any chance, the philosopher Zeno of Elea?
\zeno I am indeed, Achilles.
\achilles (scratching his head in puzzlement). Now how did he know my name?
\zeno Could I possibly persuade you two to hear me out as to why this is the case? I’ve come
      all the way to Elea from point A this afternoon, just trying to find someone who’ll pay
      some attention to my closely honed argument. But they’re all hurrying hither and thither,
      and they don’t have time. You’ve no idea how disheartening it is to meet with refusal
      after refusal. Oh, I’m sorry to burden you with my troubles, I’d just like to ask you one
      thing: Would the two of you humour a sill old philosopher for a few moments – only a
      few, I promise you – in his eccentric theories.
\achilles Oh, by all means! Please do illuminate us! I know I speak for both of us, since my
      companion, Mr. Tortoise, was only moments ago speaking of you with great veneration –
      and he mentioned especially your paradoxes.
\zeno Thank you. You see, my Master, the fifth patriarch, taught me that reality is one,
      immutable, and unchanging, all plurality, change, and motion are mere illusions of the
      sense. Some have mocked his views; but I will show the absurdity of their mockery. My
      argument is quite simple. I will illustrate it with two characters of my own Invention:
      Achilles )a Greek warrior, the fleetest of foot of all mortals), and a Tortoise. In my tale,
      they are persuaded by a passerby to run a footrace down a runway towards a distant flag
      waving in the breeze. Let us assume that, since the Tortoise is a much slowerrunner, he
      gets a head start of, say, ten rods. Now the race begins. In a few bounds Achilles has
      reached the spot where the Tortoise started.




Three-Part Invention                                                                           39
\achilles Hah!
\zeno And now the Tortoise is but a single rod ahead of Achilles. Within only a moment,
      Achilles has attained that spot.
\achilles Ho ho!
\zeno Yet, in that short moment, the Tortoise has managed to advance a slight amount. In a
      flash, Achilles covers that distance too.
\achilles Hee hee hee!
\zeno But in that very short flash, the Tortoise has managed to inch ahead by ever so little, and
      so Achilles is still behind. Now you see that in order for Achilles to catch the Tortoise,
      this game of “try-to-catch-me” will have to be played an INFINITE number of times –
      and therefore Achilles can NEVER catch up with the Tortoise.
\tortoise Heh heh heh heh!
\achilles Hmm… Hmm… Hmm… Hmm… Hmm…That argument sounds wrong to me.
      And yes, I can’t quite make out what’s wrong with it
\zeno Isn’t it a teaser? It’s my favourite paradox.
\tortoise Excuse me, Zeno, but I believe your tale illustrates the wrong principle, doe sit
      not? You have just told us what will come to known, centuries hence, as Zeno’s “Achilles
      paradox” , which shows (ahem!) that Achilles will never catch the Tortoise; but the proof
      that Motion Is Inherently Impossible (and thence that Motion Unexists) is your
      “dichotomy paradox”, isn’t that so?
\zeno Oh, shame on me. Of course, you’re right. That’s the new one about how, in going from
      A to B, one has to go halfway first – and of that stretch one also has to go halfway, and so
      on and so forth. But you see, both those paradoxes really have the same flavour. Frankly,
      I’ve only had one Great Idea – I just exploit it in different ways.
\achilles I swear, these arguments contain a flaw. I don’t quite see where, but they cannot
      be correct.
\zeno You doubt the validity of my paradox? Why not just try it out|? You see that red flag
      waving down here, at the far end of the runway?
\achilles The impossible one, based on an Escher print?
\zeno Exactly. What do you say to you and Mr. Tortoise racing for it, allowing Mr. T a fair
      head start of, well, I don’t know –
\tortoise How about ten rods?
\zeno Very good – ten rods.
\achilles Any time.
\zeno Excellent! How exciting! An empirical test of my rigorously proven Theorem! Mr.
      Tortoise, will you position yourself ten rods upwind?

(The Tortoise moves ten rods closer to the flag)

Tortoise and Achlles: Ready!
\zeno On your mark! Get set! Go!




Three-Part Invention                                                                           40



                                The MU-puzzle

                                         Formal Systems

ONE OF THE most central notions in this book is that of a formal system. The type of
formal system I use was invented by the American logician Emil Post in the 1920's, and
is often called a "Post production system". This Chapter introduces you to a formal
system and moreover, it is my hope that you will want to explore this formal system at
least a little; so to provoke your curiosity, I have posed a little puzzle.
         "Can you produce MU?" is the puzzle. To begin with, you will be supplied with a
string (which means a string of letters).* Not to keep you in suspense, that string will be
MI. Then you will be told some rules, with which you can change one string into another.
If one of those rules is applicable at some point, and you want to use it, you may, but-
there is nothing that will dictate which rule you should use, in case there are several
applicable rules. That is left up to you-and of course, that is where playing the game of
any formal system can become something of an art. The major point, which almost
doesn't need stating, is that you must not do anything which is outside the rules. We
might call this restriction the "Requirement of Formality". In the present Chapter, it
probably won't need to be stressed at all. Strange though it may sound, though, I predict
that when you play around with some of the formal systems of Chapters to come, you
will find yourself violating the Requirement of Formality over and over again, unless you
have worked with formal systems before.
         The first thing to say about our formal system-the MIU-system-is that it utilizes
only three letters of the alphabet: M, I, U. That means that the only strings of the MIU-
system are strings which are composed of those three letters. Below are some strings of
the MIU-system:

                                  MU
                                  UIM
                                  MUUMUU
                                  UIIUMIUUIMUIIUMIUUIMUIIU

* In this book, we shall employ the following conventions when we refer to strings. When the
string is in the same typeface as the text, then it will be enclosed in single or double quotes.
Punctuation which belongs to the sentence and not to the string under discussion will go outside
of the quotes, as logic dictates. For example, the first letter of this sentence is 'F', while the first
letter of 'this ‘sentence’.is 't'. When the string is in Quadrata Roman, however, quotes will
usually be left off, unless clarity demands them. For example, the first letter of Quadrata is Q.




The MU-puzzle                                                                                        41

But although all of these are legitimate strings, they are not strings which are "in your
possession". In fact, the only string in your possession so far is MI. Only by using the
rules, about to be introduced, can you enlarge your private collection. Here is the first
rule:

\rule I: If you possess a string whose last letter is I, you can add on a U at the end.\par


By the way, if up to this point you had not guessed it, a fact about the meaning of "string"
is that the letters are in a fixed order. For example, MI and IM are two different strings.
A string of symbols is not just a "bag" of symbols, in which the order doesn't make any
difference.
         Here is the second rule:

\rule II: Suppose you have Mx. Then you may add Mxx to your collection.\par


What I mean by this is shown below, in a few examples.

From MIU, you may get MIUIU.
From MUM, you may get MUMUM.
From MU, you may get MUU.

So the letter `x' in the rule simply stands for any string; but once you have decided which
string it stands for, you have to stick with your choice (until you use the rule again, at
which point you may make a new choice). Notice the third example above. It shows how,
once you possess MU, you can add another string to your collection; but you have to get
MU first! I want to add one last comment about the letter `x': it is not part of the formal
system in the same way as the three letters `M', `I', and `U' are. It is useful for us,
though, to have some way to talk in general about strings of the system, symbolically-and
that is the function of the `x': to stand for an arbitrary string. If you ever add a string
containing an 'x' to your "collection", you have done something wrong, because strings of
the MIU-system never contain "x" “s”!
Here is the third rule:

\rule III: If III occurs in one of the strings in your collection, you may make a new\par

string with U in place of III.

Examples:

     From UMIIIMU, you could make UMUMU.
     From MII11, you could make MIU (also MUI).
     From IIMII, you can't get anywhere using this rule.
       (The three I's have to be consecutive.)
     From MIII, make MU.

Don't, under any circumstances, think you can run this rule backwards, as in the
following example:



The MU-puzzle                                                                            42

From MU, make MIII               <- This is wrong.

Rules are one-way.
       Here is the final rule.

\rule IV: If UU occurs inside one of your strings, you can drop it.\par


From UUU, get U.
From MUUUIII, get MUIII.

There you have it. Now you may begin trying to make MU. Don't worry you don't get it.
Just try it out a bit-the main thing is for you to get the flavor of this MU-puzzle. Have
fun.

                                 Theorems, Axioms, Rules
The answer to the MU-puzzle appears later in the book. For now, what important is not
finding the answer, but looking for it. You probably hay made some attempts to produce
MU. In so doing, you have built up your own private collection of strings. Such strings,
producible by the rules, are called theorems. The term "theorem" has, of course, a
common usage mathematics which is quite different from this one. It means some
statement in ordinary language which has been proven to be true by a rigorous argument,
such as Zeno's Theorem about the "unexistence" of motion, c Euclid's Theorem about the
infinitude of primes. But in formal system theorems need not be thought of as statements-
they are merely strings c symbols. And instead of being proven, theorems are merely
produced, as if F machine, according to certain typographical rules. To emphasize this
important distinction in meanings for the word "theorem", I will adopt the following
convention in this book: when "theorem" is capitalized, its meaning will be the everyday
one-a Theorem is a statement in ordinary language which somebody once proved to be
true by some sort of logic argument. When uncapitalized, "theorem" will have its
technical meaning a string producible in some formal system. In these terms, the MU-
puzzle asks whether MU is a theorem of the MIU-system.

I gave you a theorem for free at the beginning, namely MI. Such "free" theorem is called
an axiom-the technical meaning again being qui different from the usual meaning. A
formal system may have zero, or several, or even infinitely many axioms. Examples of all
these types v appear in the book.

Every formal system has symbol-shunting rules, such as the four rules of the MIU-
system. These rules are called either rules of production or rules of inference. I will use
both terms.

The last term which I wish to introduce at this point is derivation. Shown below is a
derivation of the theorem MUIIU:

      (1) MI                            axiom
      (2) MII                           from (1) by rule II


The MU-puzzle                                                                           43

      (3) MIII                         from (2) by rule II
      (4) MIIIIU                       from (3) by rule I
      (5) MUIU                         from (4) by rule III
      (6) MUIUUIU                      from (5) by rule II
      (7) MUIIU                        from (6) by rule IV

A derivation of a theorem is an explicit, line-by-line demonstration of how to produce
that theorem according to the rules of the formal system. The concept of derivation is
modeled on that of proof, but a derivation is an austere cousin of a proof. It would sound
strange to say that you had proven MUIIU, but it does not sound so strange to say you
have derived MUIIU.

                           Inside and Outside the System
Most people go about the MU-puzzle by deriving a number of theorems, quite at random,
just to see what kind of thing turns up. Pretty soon, they begin to notice some properties
of the theorems they have made; that is where human intelligence enters the picture. For
instance, it was probably not obvious to you that all theorems would begin with M, until
you had tried a few. Then, the pattern emerged, and not only could you see the pattern,
but you could understand it by looking at the rules, which have the property that they
make each new theorem inherit its first letter from an earlier theorem; ultimately, then, all
theorems' first letters can be traced back to the first letter of the sole axiom MI-and that is
a proof that theorems of the MIU-system must all begin with M.
         There is something very significant about what has happened here. It shows one
difference between people and machines. It would certainly be possible-in fact it would
be very easy-to program a computer to generate theorem after theorem of the MIU-
system; and we could include in the program a command to stop only upon generating U.
You now know that a computer so programmed would never stop. And this does not
amaze you. But what if you asked a friend to try to generate U? It would not surprise you
if he came back after a while, complaining that he can't get rid of the initial M, and
therefore it is a wild goose chase. Even if a person is not very bright, he still cannot help
making some observations about what he is doing, and these observations give him good
insight into the task-insight which the computer program, as we have described it, lacks.
         Now let me be very explicit about what I meant by saying this shows a difference
between people and machines. I meant that it is possible to program a machine to do a
routine task in such a way that the machine will never notice even the most obvious facts
about what it is doing; but it is inherent in human consciousness to notice some facts
about the things one is doing. But you knew this all along. If you punch "1" into an
adding machine, and then add 1 to it, and then add 1 again, and again, and again, and
continue doing so for hours and hours, the machine will never learn to anticipate you, and
do it itself, although any person would pick up the




The MU-puzzle                                                                               44

pick up the idea, no matter how much or how well it is driven, that it i supposed to avoid
other cars and obstacles on the road; and it will never learn even the most frequently
traveled routes of its owner.
        The difference, then, is that it is possible for a machine to act unobservant; it is
impossible for a human to act unobservant. Notice I am not saying that all machines are
necessarily incapable of making sophisticated observations; just that some machines are.
Nor am I saying that all people are always making sophisticated observations; people, in
fact, are often very unobservant. But machines can be made to be totally unobservant;
any people cannot. And in fact, most machines made so far are pretty close ti being
totally unobservant. Probably for this reason, the property of being; unobservant seems to
be the characteristic feature of machines, to most people. For example, if somebody says
that some task is "mechanical", i does not mean that people are incapable of doing the
task; it implies though, that only a machine could do it over and over without eve
complaining, or feeling bored.

                             Jumping out of the System
It is an inherent property of intelligence that it can jump out of the tas which it is
performing, and survey what it has done; it is always looking for and often finding,
patterns. Now I said that an intelligence can jump out o its task, but that does not mean
that it always will. However, a little prompting will often suffice. For example, a human
being who is reading a boo may grow sleepy. Instead of continuing to read until the book
is finished he is just as likely to put the book aside and turn off the light. He ha stepped
"out of the system" and yet it seems the most natural thing in the world to us. Or, suppose
person A is watching television when person B comes in the room, and shows evident
displeasure with the situation Person A may think he understands the problem, and try to
remedy it b exiting the present system (that television program), and flipping the channel
knob, looking for a better show. Person B may have a more radio concept of what it is to
"exit the system"-namely to turn the television oft Of course, there are cases where only a
rare individual will have the vision to perceive a system which governs many peoples
lives, a system which ha never before even been recognized as a system; then such people
often devote their lives to convincing other people that the system really is there and that
it ought to be exited from!
         How well have computers been taught to jump out of the system? I w cite one
example which surprised some observers. In a computer chess: tournament not long ago
in Canada, one program-the weakest of all the competing ones-had the unusual feature of
quitting long before the game was over. It was not a very good chess player, but it at least
had the redeeming quality of being able to spot a hopeless position, and to resign then
and there, instead of waiting for the other program to go through the




The MU-puzzle                                                                            45

boring ritual of checkmating. Although it lost every game it played, it did it in style. A lot
of local chess experts were impressed. Thus, if you define "the system" as "making
moves in a chess game", it is clear that this program had a sophisticated, preprogrammed
ability to exit from the system. On the other hand, if you think of "the system" as being
"whatever the computer had been programmed to do", then there is no doubt that the
computer had no ability whatsoever to exit from that system.
        It is very important when studying formal systems to distinguish working within
the system from making statements or observations about the system. I assume that you
began the MU-puzzle, as do most people, by working within the system; and that you
then gradually started getting anxious, and this anxiety finally built up to the point where
without any need for further consideration, you exited from the system, trying to take
stock of what you had produced, and wondering why it was that you had not succeeded in
producing MU. Perhaps you found a reason why you could not produce MU; that is
thinking about the system. Perhaps you produced MIU somewhere along the way; that is
working within the system. Now I do not want to make it sound as if the two modes are
entirely incompatible; I am sure that every human being is capable to some extent of
working inside a system and simultaneously thinking about what he is doing. Actually, in
human affairs, it is often next to impossible to break things neatly up into "inside the
system" and "outside the system"; life is composed of so many interlocking and
interwoven and often inconsistent "systems" that it may seem simplistic to think of things
in those terms. But it is often important to formulate simple ideas very clearly so that one
can use them as models in thinking about more complex ideas. And that is why I am
showing you formal systems; and it is about time we went back to discussing the MIU-
system.


                             M-Mode, I-Mode, U-Mode
The MU-puzzle was stated in such a way that it encouraged some amount of exploration
within the MIU-system-deriving theorems. But it was also stated in a way so as not to
imply that staying inside the system would necessarily yield fruit. Therefore it
encouraged some oscillation between the two modes of work. One way to separate these
two modes would be to have two sheets of paper; on one sheet, you work "in your
capacity as a machine", thus filling it with nothing but M's, I's, and U's; on the second
sheet, you work "in your capacity as a thinking being", and are allowed to do whatever
your intelligence suggests-which might involve using English, sketching ideas, working
backwards, using shorthand (such as the letter `x'), compressing several steps into one,
modifying the rules of the system to see what that gives, or whatever else you might
dream up. One thing you might do is notice that the numbers 3 and 2 play an important
role, since I's are gotten rid of in three's, and U's in two's-and doubling of length (except
for the M) is allowed by rule II. So the second sheet might




The MU-puzzle                                                                              46

also have some figuring on it. We will occasionally refer back to these two modes of
dealing with a formal system, and we will call them the Mechanic mode (M-mode) and
the Intelligent mode (I-mode). To round out our mode with one for each letter of the
MIU-system, I will also mention a fin mode-the Un-mode (U-mode), which is the Zen
way of approaching thing. More about this in a few Chapters.


                                 Decision Procedures
An observation about this puzzle is that it involves rules of two opposite tendencies-the
lengthening rules and the shortening rules. Two rules (I and II) allow you to increase the
size of strings (but only in very rigid, pr scribed ways, of course); and two others allow
you to shrink strings somewhat (again in very rigid ways). There seems to be an endless
variety to the order in which these different types of rules might be applied, and this gives
hope that one way or another, MU could be produced. It might involve lengthening the
string to some gigantic size, and then extracting piece after piece until only two symbols
are left; or, worse yet, it might involve successive stages of lengthening and then
shortening and then lengthening and then shortening, and so on. But there is no guarantee
it. As a matter of fact, we already observed that U cannot be produced at all and it will
make no difference if you lengthen and shorten till kingdom come.
        Still, the case of U and the case of MU seem quite different. It is by very
superficial feature of U that we recognize the impossibility of producing it: it doesn't
begin with an M (whereas all theorems must). It is very convenient to have such a simple
way to detect nontheorems. However who says that that test will detect all nontheorems?
There may be lots strings which begin with M but are not producible. Maybe MU is one
of them. That would mean that the "first-letter test" is of limited usefulness able only to
detect a portion of the nontheorems, but missing others. B there remains the possibility of
some more elaborate test which discriminates perfectly between those strings which can
be produced by the rules and those which cannot. Here we have to face the question,
"What do mean by a test?" It may not be obvious why that question makes sense, of
important, in this context. But I will give an example of a "test" which somehow seems to
violate the spirit of the word.
        Imagine a genie who has all the time in the world, and who enjoys using it to
produce theorems of the MIU-system, in a rather methodical way. Here, for instance, is a
possible way the genie might go about it

      Step 1: Apply every applicable rule to the axiom MI. This yields two new theorems
            MIU, MII.
      Step 2: Apply every applicable rule to the theorems produced in step 1. This yields
            three new theorems: MIIU, MIUIU, MIIII.




The MU-puzzle                                                                             47

      Step 3: Apply every applicable rule to the theorems produced in step 2. This yields
      five new theorems: MIIIIU, MIIUIIU, MIUIUIUIU, MIIIIIIII, MUI.


This method produces every single theorem sooner or later, because the rules are applied
in every conceivable order. (See Fig. 11.) All of the lengthening-shortening alternations
which we mentioned above eventually get carried out. However, it is not clear how long
to wait for a given string




FIGURE 11. A systematically constructed "tree" of all the theorems of the MIU-system.
The N th level down contains those theorems whose derivations contain exactly N steps.
The encircled numbers tell which rule was employed. Is MU anywhere in this tree?


to appear on this list, since theorems are listed according to the shortness of their
derivations. This is not a very useful order, if you are interested in a specific string (such
as MU), and you don't even know if it has any derivation, much less how long that
derivation might be.
        Now we state the proposed "theoremhood-test":

      Wait until the string in question is produced; when that happens, you know it
      is a theorem-and if it never happens, you know that it is not a theorem.

This seems ridiculous, because it presupposes that we don't mind waiting around literally
an infinite length of time for our answer. This gets to the crux of the matter of what
should count as a "test". Of prime importance is a guarantee that we will get our answer
in a finite length of time. If there is a test for theoremhood, a test which does always
terminate in a finite




The MU-puzzle                                                                              48

amount of time, then that test is called a decision procedure for the given formal system.
         When you have a decision procedure, then you have a very concrete
characterization of the nature of all theorems in the system. Offhand, it might seem that
the rules and axioms of the formal system provide no less complete a characterization of
the theorems of the system than a decision procedure would. The tricky word here is
"characterization". Certainly the rules of inference and the axioms of the MIU-system do
characterize, implicitly, those strings that are theorems. Even more implicitly, they
characterize those strings that are not theorems. But implicit characterization is not
enough, for many purposes. If someone claims to have a characterization of all theorems,
but it takes him infinitely long to deduce that some particular string is not a theorem, you
would probably tend to say that there is something lacking in that characterization-it is
not quite concrete enough. And that is why discovering that a decision procedure exists is
a very important step. What the discovery means, in effect, is that you can perform a test
for theoremhood of a string, and that, even if the test is complicated, it is guaranteed to
terminate. In principle, the test is just as easy, just as mechanical, just as finite, just as full
of certitude, as checking whether the first letter of the string is M. A decision procedure
is a "litmus test" for theoremhood!
         Incidentally, one requirement on formal systems is that the set of axioms must be
characterized by a decision procedure-there must be a litmus test for axiomhood. This
ensures that there is no problem in getting off the ground at the beginning, at least. That
is the difference between the set of axioms and the set of theorems: the former always has
a decision procedure, but the latter may not.
         I am sure you will agree that when you looked at the MIU-system for the first
time, you had to face this problem exactly. The lone axiom was known, the rules of
inference were simple, so the theorems had been implicitly characterized-and yet it was
still quite unclear what the consequences of that characterization were. In particular, it
was still totally unclear whether MU is, or is not, a theorem.




The MU-puzzle                                                                                   49

                FIGURE 12. Sky Castle, by M. C.: Escher (woodcut, 1928).




The MU-puzzle                                                              50

